\scanpage{14}\mspage{13}
\uline{Définition}\quad Un \unsure{ens|ens.} ordonné $I$ est dit \uline{polyédral} si deux éléments minorés ont un inf. \deleted{(NB si $I$ est fini, cel revient à ceci : l'ens \unsure{ordonné} $I_{\leq x}$ et $I$} \deleted{\inserted{(i.e. si}} \inserted{NB cela signifie que si}
$x,y\in I$ sont tels que $I_{\leq x}\cap I_{\leq y}\neq\unsure{\phi|\varnothing}$, alors $I_{\leq x}\cap I_{\leq y}$ a un plus grand élément. Quand \deleted{\illegible{3}}\inserted{\unsure{cet ens ord|cet ens est}} est inductif, (\unsure{par ex|p. ex.} \unsure{fini}) \deleted{est \unsure{filtrant}} il revient au \unsure{même|$\widetilde{\ }$} de dire que ledit \deleted{\unsure{ens}} $I_{\leq x}\cap I_{\leq y}$ est filtrant\inserted{\unsure{croiss.}}\unsure{,)} \deleted{\unsure{donc ceci ce}} signifie
\[
\begin{aligned}
&\forall\,x,y,x',y'\in I\text{ tels que }\ x',y'\leq x,y,\quad \exists\,z\in I\text{ tel que}\\
&\hspace{8em}x',y'\leq z\leq x,y
\end{aligned}
\]
% slanted
\margin{Cela signifie aussi, en termes \unsure{topolgiques|topologiques}, que l'intersection de deux parties fermées \uline{\unsure{strictmt|strict}} irréductibles est itou -- chose des plus remarquables \unsure{certes} !}

Sous cette forme, on voit que la condition est autoduale, et signifie donc aussi que deux \unsure{élément|éléments} majorés de $I$ ont un plus grand élément.

Donc \uline{Lemme} Supposons que $\forall\,x,\deleted{\unsure{y,x',y'}}x'\in I$ tels que $x'\leq x$, $I^{\mathrm{opp}}_{x'\leq\,\cdot\,\leq x}$ soit inductif. \deleted{Alors} (\unsure{p.ex.|p. ex.} fini). Alors si $I$ est polyédral, $I^{\mathrm{opp}}$ l'est, et l'inverse est vrai si les $I_{x'\leq\,\cdot\,\leq x}$ sont inductifs. \unsure{?}\unsure{à|--}.

Si \deleted{\unsure{Soit}} $\mathcal{A}=(\mathcal{F},\unsure{\leq^{\prime}|\leq},\ll)$ est un atelier, $F\in\unsure{\mathcal{J}|\mathcal{F}}$, la \uline{figure} \uline{\unsure{ou|$\sigma$}} $F$ est dite \uline{polyédrale} si $\widetilde{F}$ est \unsure{polyédral|polyédrale}, i.e. si deux strates de $F$ non disjointes ont comme intersection une strate de $F$. L'atelier $\mathcal A$ est dit \uline{polyédral} si \deleted{\unsure{tout}} figure de $\mathcal A$ est polyédral, i.e si on a ceci (en termes de $\unsure{M|\mathcal{M}}$) :

% slanted
\margin{(définition des ateliers polyédraux, ou condition de la stricte \unsure{intersection}}
\uline{At} \uline{pol} 1\quad $\forall\,X,Y\in M$, si $X\underset{M}{\mathrel{>\mkern-12mu<}}Y$ et $X$ et $Y$ non disjoints (i.e. $\{X,Y\}$ minoré dans $M$) alors \deleted{\unsure{on}} $\mathrm{Inf}^{M,\leq}(X,Y)$ existe.

Un ens. ordonné \inserted{$I$} est dit \uline{simplicial}, s'il est polyédral, et si pour tout $x\in I$, $S=I_{\leq x}$ est un \uline{simplexe combinatoire}, i.e. isomorphe à \deleted{\unsure{l'}}\inserted{un} ensemble ordonné ${\unsure{\mathcal{P}|\mathscr{P}}}^*(J)$ \inserted{($\mathcal P^*\unsure{[J]|(J)}={}$ ens. des parties \uline{non vides} de $J$)}, où $J$ est un ens \uline{fini}.

Pour un ens. ordonné $S$, il revient au $=$ de dire que $S$ est un simplexe combinatoire, ou que $S$ est \inserted{polyédral} fini, qu'il a un plus grand élément, et que chaque élément non minimal a au moins deux prédécesseurs. [Alors on voit par induction sur $\mathrm{card}\,S$ que \unsure{tt|tout} élément de $S$ est
